% Standalone: architektura kodeka MLX (train_mlx) — do wklejenia do pracy lub kompilacji osobno
\documentclass[11pt,a4paper]{article}

\usepackage{fontspec}
\usepackage{polyglossia}
\setdefaultlanguage{polish}
\defaultfontfeatures{Ligatures=TeX}
\setmainfont{Times New Roman}
\setsansfont{Arial}
\setmonofont{Menlo}
\usepackage{amsmath,amssymb}
\usepackage{microtype}
\usepackage{geometry}
\geometry{margin=2.5cm}
\usepackage{hyperref}
\hypersetup{colorlinks=true, linkcolor=blue, urlcolor=blue, citecolor=blue}

\title{Architektura neuronalnego kodeka audio\\w implementacji MLX}
\author{}
\date{}

\begin{document}
\maketitle

\begin{abstract}
Opis dotyczy modelu \texttt{MLXCodec} z pliku \texttt{tools/train\_mlx.py}: enkoder splotowy, kwantyzacja wektorowa resztowa (RVQ) oraz dekoder odbijający strukturę enkodera, uczony na MLX z funkcją kosztu łączącej błąd czasowy, wieloskalową stratę STFT oraz termy związane z różnorodnością kodów.
\end{abstract}

\subsection{Przepływ sygnału}

Niech mono próbki mają kształt $x \in \mathbb{R}^{B \times T \times 1}$ przy częstotliwości próbkowania $f_s = 16\,\mathrm{kHz}$.
Model realizuje odwzorowanie
\begin{equation}
  x \;\mapsto\; z = f_{\mathrm{enc}}(x) \;\mapsto\; \hat z = \mathrm{RVQ}(z) \;\mapsto\; \hat x = g_{\mathrm{dec}}(\hat z),
  \label{eq:pipeline}
\end{equation}
gdzie $z$ oraz skwantowane $\hat z$ mają wymiar ukryty $D = 128$ wzdłuż skompresowanej osi czasu. Architektura jest spójna z linią prac typu SoundStream / EnCodec \cite{soundstream,encodec} oraz enkoderów złożonych z splotów ze stride (SEANet).

\subsection{Enkoder}

Enkoder składa się z $L = 8$ stopni ze stride czasowym $2$; każdy stopień używa splotów jednowymiarowych z jądrem $7$, nieliniowości \textsc{GELU} oraz rosnącej liczby kanałów
\begin{equation}
  (c_0, \ldots, c_{L}) = (1,\, 24,\, 32,\, 48,\, 64,\, 96,\, 128,\, 192,\, 256).
\end{equation}
Przed każdym zmniejszeniem rozdzielczości czasu o połowę można wstawić $n_1 \ge 0$ bloków stride-$1$ (domyślnie $n_1 = 1$) przy szerokości wejściowej danego stopnia, co daje płynniejszą refinemację cech. Łączny stride czasowy wynosi
\begin{equation}
  S = 2^L = 256,
\end{equation}
więc częstotliwość ramek latentnych to $f_s / S$. Wyjście enkodera to splot do wymiaru $D = 128$.

\subsection{Kwantyzacja resztowa (RVQ)}

Stosuje się $N_q = 2$ etapów RVQ \cite{soundstream,encodec}: w etapie $i$ kwantyzowany jest resztkowy wektor $r^{(i)} = z - \sum_{j<i} \hat z^{(j)}$.
Każdy etap ma osobny kodebook $\mathcal{E}^{(i)} \subset \mathbb{R}^D$ o $K = 128$ wektorach.

\paragraph{Tryb kosinusowy.}
Włączony jest wariant \emph{cosine VQ}: reszta i wektory kodów są normalizowane do jednostkowych kierunków; indeks wybiera minimalną kwadratową odległość kątową na sferze. Rekonstrukcja etapu skaluje wybrany kierunek normą bieżącej reszty; gradient przechodzi przez resztę i wybrane kody mechanizmem STE (straight-through estimator).

Strata kwantyzacji na etap to klasyczna suma commitmentu (waga $\beta$) oraz człon zbliżający kody do zatrzymanego gradientu reszty; suma po etapach jest mnożona przez $\lambda_{\mathrm{vq}}$.

Nominalnie na ramkę latentną przypada $N_q \log_2 K$ bitów informacji dyskretnej; przy powyższych ustawieniach jest to $14$ bitów na ramkę, co przy $f_s = 16\,\mathrm{kHz}$ i $S = 256$ daje szacunkową przepływność rzędu $0{,}9\,\mathrm{kbps}$ (tylko warstwa dyskretna, bez nagłówków kontenera).

\subsection{Dekoder}

Dekoder jest lustrzanym odbiciem enkodera: \texttt{ConvTranspose1d} ze stride $2$, opcjonalne bloki stride-$1$ po zwiększeniu rozdzielczości, końcowy splot do jednego kanału. Długość czasu wyrównuje się do wejścia; na wyjściu stosuje się $\tanh$.

\subsection{Funkcja kosztu}

Niech $\lVert \cdot \rVert_1$ oznacza średni błąd $\ell_1$ po próbkach i batchu. Strata wieloskalowa STFT to średnia po skalach $(n_{\mathrm{fft}}, \mathrm{hop})$ z $\ell_1$ na log-modułach transformacji (okno Hanna). Efektywna waga STFT rośnie liniowo w czasie treningu od ułamka początkowego do pełnej wartości $\lambda_{\mathrm{stft}}$ (ramp), aby wczesna faza nie zdominowała uczenia RVQ.

Dla różnorodności kodów używa się entropii rozkładu marginalnego $\bar p$ po batchu: z macierzy odległości do kodów buduje się $\mathrm{softmax}(-d/\tau)$, marginalizuje po próbkach i czasie, oblicza się entropię Shannona $H(\bar p)$ oraz dodaje się karę
\begin{equation}
  \lambda_m \max\bigl(0,\, \log K - H(\bar p)\bigr),
  \label{eq:marginal}
\end{equation}
przy czym na początku treningu $\lambda_m$ może być przez stały czas zwiększona (\emph{boost}). Opcjonalnie dodaje się kary kosinusowej zgodności fali lub reset martwych wektorów kodu przy wykrytym kolapsie.

Całkowity koszt jest sumą składników czasowych, STFT (ważonej), RVQ oraz \eqref{eq:marginal} (i opcjonalnych dodatków). Optymalizacja: Adam, globalne przycinanie normy gradientu.

\subsection{Dane treningowe}

W trybie z korpusem mówcy wczytywane są segmenty z drzewa plików audio (\texttt{.flac}, \texttt{.wav}); amplituda normalizowana jest per próbka, domyślna długość segmentu to $16384$ próbek ($\approx 1\,\mathrm{s}$ przy $16\,\mathrm{kHz}$).

\begin{thebibliography}{9}

\bibitem{soundstream}
N.~Zeghidour, A.~d'Autume, J.~Kopf, J.~Zeghidour, N.~Usunier.
\newblock \emph{SoundStream: An End-to-End Neural Audio Codec.}
\newblock IEEE/ACM TASLP, 29:2858--2874, 2021.

\bibitem{encodec}
A.~Défossez, J.~Copet, G.~Synnaeve, Y.~Adi.
\newblock \emph{High Fidelity Neural Audio Compression.}
\newblock arXiv:2210.13438, 2022.

\end{thebibliography}

\end{document}
